% ============================================================================
% preamble.tex — The CAEM Book
% Defines the visual language: pedagogical callout boxes, code styling,
% theorem environments, colour palette, and typography.
% ============================================================================

% --- Page geometry & typography ---------------------------------------------
\usepackage[a4paper,top=1in,bottom=1in,left=1.1in,right=1.1in]{geometry}
\usepackage[T1]{fontenc}
\usepackage[utf8]{inputenc}
\usepackage{lmodern}
\usepackage{microtype}
\usepackage{setspace}
\onehalfspacing

% --- Maths ------------------------------------------------------------------
\usepackage{amsmath}
\usepackage{amssymb}
\usepackage{amsthm}
\usepackage{bm}
\usepackage{mathtools}

% --- Tables, graphics, colour -----------------------------------------------
\usepackage{graphicx}
\usepackage{booktabs}
% --- Float placement: keep figures/tables in the section that references them
\usepackage{float}                  % [H] exact placement for small reference tables
\usepackage[section]{placeins}      % \FloatBarrier before every \section (no cross-section drift)
% Allow more floats to sit on a text page near their reference (less deferral/drift)
\setcounter{topnumber}{4}
\setcounter{bottomnumber}{3}
\setcounter{totalnumber}{6}
\renewcommand{\topfraction}{0.92}
\renewcommand{\bottomfraction}{0.85}
\renewcommand{\textfraction}{0.06}
\renewcommand{\floatpagefraction}{0.65}
\usepackage{array}
\usepackage{multirow}
\usepackage{tabularx}
\usepackage[table,dvipsnames]{xcolor}
\usepackage{enumitem}
\usepackage{caption}

% --- Diagrams & plots -------------------------------------------------------
\usepackage{tikz}
\usetikzlibrary{arrows.meta,positioning,shapes.geometric,fit,backgrounds,
                calc,decorations.pathreplacing,shadows.blur,patterns}
\usepackage{pgfplots}
\pgfplotsset{compat=1.18}

% --- Code listings ----------------------------------------------------------
\usepackage{listings}

% --- Callout boxes ----------------------------------------------------------
\usepackage[most]{tcolorbox}
\tcbuselibrary{skins,breakable}

% --- Bibliography (shares the thesis reference database) --------------------
\usepackage[backend=biber,style=numeric-comp,sorting=none,maxbibnames=99]{biblatex}
\addbibresource{../thesis_report/bibliography/references.bib}

% --- Cross-refs & links -----------------------------------------------------
\usepackage[hidelinks,colorlinks=true,linkcolor=NavyBlue,
            citecolor=NavyBlue,urlcolor=Maroon]{hyperref}
\usepackage[noabbrev,capitalize]{cleveref}

% --- Section / chapter styling ----------------------------------------------
\usepackage{titlesec}
\usepackage{titletoc}

% Shorter, non-uppercase running heads (the default "CHAPTER N. FULL TITLE" in caps
% overflows the header band on long chapter titles). Mixed-case, no "Chapter" word.
\renewcommand{\chaptermark}[1]{\markboth{\thechapter.\ #1}{}}
\renewcommand{\sectionmark}[1]{\markright{\thesection.\ #1}}

% ============================================================================
% COLOUR PALETTE  (a calm, readable teaching palette)
% ============================================================================
\definecolor{cbInk}{HTML}{1A1A2E}        % body ink
\definecolor{cbIntuition}{HTML}{1565C0}  % blue  — intuition
\definecolor{cbFormula}{HTML}{6A1B9A}    % purple — precise/maths
\definecolor{cbExample}{HTML}{2E7D32}    % green — worked example
\definecolor{cbScenario}{HTML}{00838F}   % teal  — scenario walkthrough
\definecolor{cbPitfall}{HTML}{C62828}    % red   — common pitfall
\definecolor{cbCode}{HTML}{37474F}       % slate — code
\definecolor{cbDefense}{HTML}{5D4037}    % brown — examiner Q&A
\definecolor{cbNumbers}{HTML}{EF6C00}    % orange— by the numbers
\definecolor{cbCodeBg}{HTML}{F5F7FA}     % code background
\definecolor{cbCodeKw}{HTML}{0B5FA5}     % code keywords
\definecolor{cbCodeStr}{HTML}{2E7D32}    % code strings
\definecolor{cbCodeCom}{HTML}{8A8A8A}    % code comments

% ============================================================================
% CODE STYLE  (Python excerpts straight from the CAEM source)
% ============================================================================
\lstdefinestyle{caempy}{
  language=Python,
  basicstyle=\ttfamily\footnotesize\color{cbInk},
  keywordstyle=\color{cbCodeKw}\bfseries,
  stringstyle=\color{cbCodeStr},
  commentstyle=\color{cbCodeCom}\itshape,
  numberstyle=\tiny\color{cbCodeCom},
  numbers=left,
  numbersep=8pt,
  backgroundcolor=\color{cbCodeBg},
  frame=single,
  framerule=0pt,
  rulecolor=\color{cbCodeBg},
  breaklines=true,
  breakatwhitespace=true,
  showstringspaces=false,
  tabsize=4,
  xleftmargin=10pt,
  xrightmargin=4pt,
  aboveskip=1ex,
  belowskip=1ex,
  literate={τ}{{$\tau$}}1 {α}{{$\alpha$}}1 {θ}{{$\theta$}}1 {ρ}{{$\rho$}}1
           {λ}{{$\lambda$}}1 {φ}{{$\phi$}}1 {≥}{{$\geq$}}1 {≤}{{$\leq$}}1
           {×}{{$\times$}}1 {→}{{$\rightarrow$}}1,
}
\lstset{style=caempy}

% ============================================================================
% PEDAGOGICAL CALLOUT BOXES
% ============================================================================
% Generic builder: a titled, coloured, breakable box with an icon-ish tag.
\newtcolorbox{cbbox}[3]{%
  enhanced, breakable, drop fuzzy shadow=black!18, colback=#1!4, colframe=#1,
  coltitle=white, fonttitle=\bfseries\sffamily,
  boxrule=0.8pt, arc=2pt, left=6pt, right=6pt, top=4pt, bottom=4pt,
  title={#2}, label={#3},
  attach boxed title to top left={yshift=-2mm, xshift=4mm},
  boxed title style={colback=#1, arc=1pt, boxrule=0pt},
}

% Intuition — the plain-English mental model
\newtcolorbox{intuition}[1][Intuition]{enhanced, breakable, drop fuzzy shadow=black!18, colback=cbIntuition!4,
  colframe=cbIntuition, coltitle=white, fonttitle=\bfseries\sffamily,
  boxrule=0.8pt, arc=2pt, title={\,#1},
  attach boxed title to top left={yshift=-2.4mm,xshift=4mm},
  boxed title style={colback=cbIntuition,arc=1pt,boxrule=0pt}}

% Precise — the exact maths / mechanism
\newtcolorbox{precise}[1][The precise version]{enhanced, breakable, drop fuzzy shadow=black!18, colback=cbFormula!4,
  colframe=cbFormula, coltitle=white, fonttitle=\bfseries\sffamily,
  boxrule=0.8pt, arc=2pt, title={\,#1},
  attach boxed title to top left={yshift=-2.4mm,xshift=4mm},
  boxed title style={colback=cbFormula,arc=1pt,boxrule=0pt}}

% Worked example
\newtcolorbox{workedexample}[1][Worked example]{enhanced, breakable, drop fuzzy shadow=black!18,
  colback=cbExample!4, colframe=cbExample, coltitle=white,
  fonttitle=\bfseries\sffamily, boxrule=0.8pt, arc=2pt, title={\,#1},
  attach boxed title to top left={yshift=-2.4mm,xshift=4mm},
  boxed title style={colback=cbExample,arc=1pt,boxrule=0pt}}

% Scenario — a query traced through the system
\newtcolorbox{scenario}[1][Scenario]{enhanced, breakable, drop fuzzy shadow=black!18,
  colback=cbScenario!4, colframe=cbScenario, coltitle=white,
  fonttitle=\bfseries\sffamily, boxrule=0.8pt, arc=2pt, title={\,#1},
  attach boxed title to top left={yshift=-2.4mm,xshift=4mm},
  boxed title style={colback=cbScenario,arc=1pt,boxrule=0pt}}

% Common pitfall / trap
\newtcolorbox{pitfall}[1][Common pitfall]{enhanced, breakable, drop fuzzy shadow=black!18,
  colback=cbPitfall!4, colframe=cbPitfall, coltitle=white,
  fonttitle=\bfseries\sffamily, boxrule=0.8pt, arc=2pt, title={\,#1},
  attach boxed title to top left={yshift=-2.4mm,xshift=4mm},
  boxed title style={colback=cbPitfall,arc=1pt,boxrule=0pt}}

% From the code — a pointer to the actual implementation
\newtcolorbox{fromcode}[1][From the code]{enhanced, breakable, drop fuzzy shadow=black!18,
  colback=cbCode!4, colframe=cbCode, coltitle=white,
  fonttitle=\bfseries\sffamily, boxrule=0.8pt, arc=2pt, title={\,#1},
  attach boxed title to top left={yshift=-2.4mm,xshift=4mm},
  boxed title style={colback=cbCode,arc=1pt,boxrule=0pt}}

% If an examiner asks — defence framing
\newtcolorbox{defence}[1][If an examiner asks]{enhanced, breakable, drop fuzzy shadow=black!18,
  colback=cbDefense!4, colframe=cbDefense, coltitle=white,
  fonttitle=\bfseries\sffamily, boxrule=0.8pt, arc=2pt, title={\,#1},
  attach boxed title to top left={yshift=-2.4mm,xshift=4mm},
  boxed title style={colback=cbDefense,arc=1pt,boxrule=0pt}}

% By the numbers — the registered constants
\newtcolorbox{bythenumbers}[1][By the numbers]{enhanced, breakable, drop fuzzy shadow=black!18,
  colback=cbNumbers!5, colframe=cbNumbers, coltitle=white,
  fonttitle=\bfseries\sffamily, boxrule=0.8pt, arc=2pt, title={\,#1},
  attach boxed title to top left={yshift=-2.4mm,xshift=4mm},
  boxed title style={colback=cbNumbers,arc=1pt,boxrule=0pt}}

% ============================================================================
% THEOREM-LIKE ENVIRONMENTS (used in the theory chapter)
% ============================================================================
\theoremstyle{plain}
\newtheorem{theorem}{Theorem}[chapter]
\newtheorem{corollary}[theorem]{Corollary}
\newtheorem{lemma}[theorem]{Lemma}
\theoremstyle{definition}
\newtheorem{definition}[theorem]{Definition}
\theoremstyle{remark}
\newtheorem{remark}[theorem]{Remark}

% ============================================================================
% CONVENIENCE MACROS  (CAEM symbols used throughout)
% ============================================================================
\newcommand{\upre}{u_{\text{pre}}}
\newcommand{\ustored}{u_{\text{stored}}}
\newcommand{\utoken}{u_{\text{token}}}
\newcommand{\udrop}{u_{\text{dropout}}}
\newcommand{\uint}{u_{\text{internal}}}
\newcommand{\savg}{s_{\text{avg}}}
\newcommand{\pentail}{p_{\text{entail}}}
\newcommand{\pgmax}{p_{\text{ground,max}}}
\newcommand{\pgmean}{p_{\text{ground,mean}}}
\newcommand{\pgatom}{p_{\text{ground,atomic}}}
\newcommand{\qarel}{q_{\text{a,rel}}}
\newcommand{\pik}{p_{\text{IK}}}
\newcommand{\prag}{p_{\text{RAG}}}
\newcommand{\taustore}{\tau_{\text{store}}}
\newcommand{\taudefer}{\tau_{\text{defer}}}
\newcommand{\tautrain}{\tau_{\text{train}}}
\newcommand{\tauretro}{\tau_{\text{retro}}}
\newcommand{\tauruc}{\tau_{\text{RUC}}}
\newcommand{\phisafe}{\phi_{\text{safe}}}
\newcommand{\rhomin}{\rho_{\min}}
\newcommand{\cconv}{c_{\text{conv}}}
\newcommand{\caem}{\textsc{CAEM}}

% A small inline "code path" style for naming files/functions in prose.
\newcommand{\code}[1]{\textcolor{cbCode}{\texttt{\small #1}}}

% ============================================================================
% HEADER STYLING  (coloured chapter + section headers, a calm teaching look)
% ============================================================================
\definecolor{cbHeader}{HTML}{14213D}   % deep navy for chapter numerals

\titleformat{\chapter}[display]
  {\sffamily\bfseries}
  {\raggedright\large\sffamily\color{cbIntuition}%
     \MakeUppercase{\chaptertitlename}\ \thechapter}
  {6pt}
  {\Huge\color{cbHeader}\raggedright}
  [\vspace{4pt}{\color{cbIntuition}\titlerule[1.2pt]}]
\titlespacing*{\chapter}{0pt}{6pt}{22pt}

\titleformat{\section}
  {\sffamily\Large\bfseries\color{cbHeader}}
  {\color{cbIntuition}\thesection}{0.7em}{}
\titleformat{\subsection}
  {\sffamily\large\bfseries\color{cbHeader}}
  {\color{cbIntuition}\thesubsection}{0.6em}{}

% Part pages: large centred coloured numeral + title
\titleformat{\part}[display]
  {\centering\sffamily\bfseries}
  {\Large\color{cbIntuition}\MakeUppercase{\partname}\ \thepart}
  {10pt}
  {\Huge\color{cbHeader}}
  [\vspace{6pt}{\color{cbIntuition}\titlerule[1.5pt]}]
